\section{Related Work}
\subsection{Prompt injection and automated attack generation}
Prompt injection attacks have become a major security concern for LLM-integrated systems 
because trusted instructions and untrusted content are often processed within the same model 
context \cite{greshake2023indirect,liu2024formalizing}. 
Greshake et al. \cite{greshake2023indirect} showed that attackers can embed malicious 
instructions in external content, such as web pages or documents, which are later executed when consumed by an LLM application. 
Subsequent work formalized prompt injection attacks and introduced systematic frameworks 
for evaluating attacks and defenses across models and tasks \cite{liu2024formalizing}. 
Beyond explicit injections, PromptRobust demonstrated that even small perturbations at the 
character, word, sentence, or semantic level can significantly affect model behavior \cite{zhu2023promptrobust}. 
Together, these findings highlight that attack success depends on multiple factors, including the model, task, prompt structure, and location of attacker-controlled input.

Research subsequently shifted from hand-crafted prompt injections toward automated attack generation, 
motivated by the brittleness and limited transferability of manually designed prompts \cite{zou2023universal,shi2024judgedeceiver}. 
Universal adversarial suffixes demonstrated that optimization-based search can automatically generate attacks that transfer across aligned models, 
including black-box systems \cite{zou2023universal}. 
This idea was later extended to LLM-as-a-Judge settings, where JudgeDeceiver searches for adversarial token sequences 
that cause the judge to prefer a target response \cite{shi2024judgedeceiver}. 
These studies show that effective attacks can be generated automatically rather than manually crafted, making malicious behavior increasingly difficult to identify. 
As a result, attacks may appear as benign retrieved content, subtle prompt perturbations, or optimized transferable suffixes rather than explicit suspicious instructions \cite{greshake2023indirect,zhu2023promptrobust,zou2023universal}.


\subsection{Defenses for LLM-integrated systems}
Existing defenses for LLM-integrated systems are commonly divided into soft and hard defenses \cite{liu2024formalizing,yao2024survey}. 
This taxonomy captures the distinction between mitigation techniques applied at inference time and defenses that modify the underlying system design. 
% We review representative approaches from both categories below.

\paragraph{Soft defenses.}
Soft defenses aim to mitigate prompt injection without modifying the underlying model or application architecture \cite{liu2024formalizing}. 
Liu et al. \cite{liu2024formalizing} benchmark a diverse set of soft defenses, ranging from prompt modifications 
to inference-time detection techniques. 
For indirect prompt injection, BIPIA introduces boundary-awareness mechanisms and explicit reminders that 
encourage the model to treat retrieved content as data rather than executable instructions \cite{yi2025benchmarking}.
These methods are practical and easy to deploy, but their effectiveness often decreases against adaptive attacks or when models fail to consistently follow prompt-level safeguards \cite{liu2024formalizing,yi2025benchmarking}.

\paragraph{Hard defenses.}
In contrast, hard defenses seek to address prompt injection at the architectural level by 
enforcing stronger trust boundaries, context separation, or input authenticity guarantees \cite{chen2024struq,jiang2023identifying}. 
The StruQ framework separates system instructions and user data into structured channels and fine-tunes the model to follow only the instruction channel \cite{chen2024struq}. 
By construction, this approach treats prompt injection as a context-separation problem rather than solely a malicious-text detection problem. 
Shield adopts a broader system perspective by defining integrity, source identification, attack detectability
and utility preservation as required safety properties for LLM-integrated applications \cite{jiang2023identifying}. 
Survey studies place these mechanisms within a wider defense toolbox that also includes corpus cleaning, robustness optimization and safe instruction tuning \cite{yao2024survey}. 

Overall, recent work suggests a shift from prompt-level mitigation toward architectural defenses that provide stronger guarantees through context separation, authenticated data flows and safer-by-design interfaces \cite{yi2025benchmarking,chen2024struq,jiang2023identifying}.


\subsection{Threat modeling and automation for LLM systems}
LLM-based threat modeling systems are emerging as tools that adapt classical
security methods to AI-integrated applications \cite{tete2024threat}. 
Threat modeling for LLM-powered applications combines STRIDE, Shostack's four-question framework and DREAD to
reason about risks such as prompt injection, data poisoning, jailbreaking, and information disclosure. 
Automated LLM integration frameworks extend this idea by using NLP,
machine learning, behavioral profiling, and anomaly detection to identify and assess LLM-specific 
threats\cite{gadade2025framework}. 
STRIDE GPT illustrates the practical direction of this domain by using LLMs to generate threat models, attack trees, mitigations,
and security test cases from system descriptions or code context
\cite{adams2024stridegpt}.

This automation builds on earlier work that turns threat models into executable
or queryable security artifacts \cite{xu2012automated}. 
Formal threat models can be transformed into attack paths and runnable security tests, making the model itself a basis for
systematic validation. 
Security design flaws can also be detected automatically by representing designs as annotated Data Flow Diagrams
and applying model query patterns derived from inspection guidelines \cite{tuma2020automating}. 
The Meta Attack Language supports reusable formal attack descriptions by defining domain-specific languages for threat modeling
and attack simulation \cite{johnson2018meta}.
